% !TEX root = ../../main.tex
\subsection{Exercise 1}
\task{Requirement}
The main memory has a 32-bit address space, the cache size is 4MB, one block is
256B and the access unit of the system is 1 byte. Determine the tag, index and
byte offset for the following cache configurations:
(a) Direct mapped, (b) 4-way set associative, (c) Fully associative.

\task{Solution}
\begin{itemize}[itemsep=2pt]
  \item Address width: 32 bits.
  \item Byte offset: $256\text{B} / 1\text{B} = 256 = 2^{8}$ $\Rightarrow$ 8 bits.
  \item Number of blocks in cache: $4\text{MB} / 256\text{B} = 2^{22}/2^{8} = 2^{14}$ blocks.
  \item Number of sets $=$ number of blocks / number of ways; Index $= \log_2(\text{sets})$;
        Tag $= 32 - \text{Index} - \text{Offset}$.
\end{itemize}

\begin{table}[H]
  \centering
  \begin{tabular}{lcccc}
    \toprule
    \textbf{Configuration} & \textbf{Sets} & \textbf{Tag} & \textbf{Index} & \textbf{Byte offset} \\
    \midrule
    (a) Direct mapped           & $2^{14}$ & 10 & 14 & 8 \\
    (b) 4-way set associative   & $2^{12}$ & 12 & 12 & 8 \\
    (c) Fully associative       & $1$      & 24 & 0  & 8 \\
    \bottomrule
  \end{tabular}
  \caption{Exercise 1 -- Address fields (bits)}
\end{table}

\task{Implementation}
The program stores the parameters in the data segment, computes the offset as
$\log_2(\text{block}/\text{unit})$, the number of sets for each configuration
(ways $= 1, 4$, and ways $=$ number of blocks for fully associative), then the
index and tag bits, and prints them.
\lstinputlisting[caption={Exercise1.asm}]{Exercise1.asm}

\task{Output}
\screenshot{images/ex1.png}{Exercise 1 -- Output in MARS}
